% The paper's plaintext trap example, specialized to c=0.
\begin{frame}{MBQC: Trap Copies}
  \centering
  \setbeamercolor{itemize item}{fg=CVQCTeal}
  \begin{tikzpicture}
    \path[use as bounding box] (-6.6,-3.15) rectangle (6.6,3.05);

    \node[anchor=north west,inner sep=0pt,text=CVQCInk] at (-6.0,2.65) {%
      \begin{minipage}{12cm}
        \raggedright
        \begin{itemize}
          \setlength{\itemsep}{0.65em}
          \item<1-> Use traps to catch a cheating prover.
          \item<2-> Our construction relies on the trap idea from
            \href{https://doi.org/10.1103/PhysRevA.96.012303}{[FK17]},
            but instead places the tests in separate copies.
        \end{itemize}
      \end{minipage}%
    };

    \only<3>{
      \node[inner sep=0pt,text=CVQCInk] at (0,-1.05) {%
        \scalebox{1.8}{%
          $\Qcircuit @C=1em @R=1em {
            \lstick{\ket{+}} & \ctrl{1} & \gate{H} & \meter & \rstick{m=0}\cw \\
            \lstick{\ket{0}} & \control \qw & \qw & \qw & \rstick{\ket{0}}\qw
          }$%
        }%
      };
    }
  \end{tikzpicture}

  \note<1>{
    \footnotesize
    The preceding plaintext circuit explains honest computation, but a
    malicious prover may perform a different computation or lie about a
    measurement result. Introduce tests with predictable honest outcomes.
    Such traps let the verifier check the prover's reports. This is the
    motivation, not a claim that one trap detects every possible deviation.
  }
  \note<2>{
    \footnotesize
    The second bullet is copied from the manuscript's Technical Overview,
    subsection Trap Copy and Cut-and-Choose. FK17 is Fitzsimons--Kashefi,
    Unconditionally Verifiable Blind Quantum Computation (2017). Our
    construction places the tests in separate MBQC copies, using the same
    public graph and schedule as real copies. The verifier secretly chooses
    which copies to test. A trap copy uses trap and dummy vertices, with
    deterministic decoded outcomes at the checked measurements. We later
    use hidden trap copies and cut-and-choose to control the number of
    corrupted real copies; some errors can survive these tests.
  }
  \note<3>{
    \footnotesize
    This is the paper's Figure ``A plaintext trap,'' specialized to c=0.
    In the general example the upper input is $T^{-c}\ket{+}$ and the lower
    input is $\ket{0}$. Setting c=0 gives the displayed plus and zero inputs,
    and the gate $HT^0$ is H. CZ does nothing because the lower input is zero.
    H maps the upper plus state to zero, so its measurement deterministically
    returns m=0. The lower wire stays in zero. A reported result of one
    exposes a deviation. This is a test circuit with different inputs from
    the computation gadget: the lower output formula from the real example
    assumed a fresh plus state and should not be substituted here.
    The plaintext example alone does not hide the trap. Next, one-time pad
    the inputs and mask the measurement instructions, so the prover cannot
    distinguish real and trap copies and selectively attack only the real
    ones. The full trap-copy protocol uses plaintext exponent zero.
  }
\end{frame}
